% GENERATED FILE. Edit canonical lessons, then run npm run generate:books.
% canonical-source-hash: fnv1a64:e7d30c226ee5b378
% canonical-lessons: JA-C383-shimaguni, JA-C383-wan, JA-C383-kaikyou, JA-C383-kaisuiyoku, JA-C383-sukitouru, JA-R383-review-378-383

\chapter{An Island Country}
\label{ch:ja-an-island-country}
\hlchaptermodality{383}

\begin{chapteropening}
\textbf{By the end of this chapter:} \emph{I can talk about bays, straits and swimming in the sea.}

Complete the last lesson of chapter 383: I can talk about bays, straits and swimming in the sea.
\end{chapteropening}

\section[\texorpdfstring{\ja{しまぐに} (shimaguni)}{しまぐに (shimaguni)}]{\ja{しまぐに} (shimaguni) — an island country}

\label{lesson:JA-C383-shimaguni}



\begin{quote}
Before the new one: say the Japanese for the horizon (over the sea), then the Japanese for to sparkle.
\end{quote}

\subsection*{You'll want to know: \ja{しまぐに}}
\textbf{\ja{しまぐに}} — \emph{shimaguni} — \textquotedblleft{}an island country\textquotedblright{}.

\begin{quote}
\textbf{\ja{日本}\ja{は} \ja{しまぐにです。}} — \emph{nihon wa shimaguni desu} — \textquotedblleft{}Japan is an island country.\textquotedblright{}
\end{quote}

\begin{cousinweb}
\textbf{\ja{しまぐに}} (\emph{shimaguni}) is a native Japanese word (wago), inherited rather than borrowed from Chinese, built from \emph{shima}, \textquotedblleft{}island\textquotedblright{} and \emph{kuni}, \textquotedblleft{}country\textquotedblright{}, joined as \emph{shimaguni}; its kanji were attached later as kun'yomi (JMdict).
\end{cousinweb}

\begin{grammarlens}[title={What you've built}]
One more word for getting around.
\end{grammarlens}

\subsection*{Your turn}
\begin{itemize}
\raggedright
  \item \emph{Say it:} \emph{shimaguni}
  \item \emph{Say it:} \emph{shimaguni}, once more
  \item \emph{Say it:} \emph{nihon wa shimaguni desu} — Japan is an island country.
  \item \emph{Recall:} say the Japanese for the horizon (over the sea), then the Japanese for to sparkle, then say \emph{shimaguni} again
\end{itemize}

\begin{tcolorbox}[breakable,colback=teal!4,colframe=teal!35!black,title={Before you move on}]
What does \ja{しまぐに} mean? (\textquotedblleft{}An island country\textquotedblright{}.) Say it once more.
\end{tcolorbox}
\section[\texorpdfstring{\ja{わん} (wan)}{わん (wan)}]{\ja{わん} (wan) — a bay}

\label{lesson:JA-C383-wan}



\begin{quote}
Before the new one: say the Japanese for to sparkle, then the Japanese for an island country.
\end{quote}

\subsection*{You'll want to know: \ja{わん}}
\textbf{\ja{わん}} — \emph{wan} — \textquotedblleft{}a bay\textquotedblright{}.

\begin{quote}
\textbf{\ja{わんの} \ja{なかは} \ja{なみが} \ja{しずかです。}} — \emph{wan no naka wa nami ga shizuka desu} — \textquotedblleft{}Inside the bay the waves are calm.\textquotedblright{}
\end{quote}

\begin{cousinweb}
\textbf{\ja{わん}} (\emph{wan}) is Sino-Japanese (kango): one Chinese character read with its on'yomi, a reading borrowed from Chinese: \emph{wan}, \textquotedblleft{}bay\textquotedblright{} (Middle Chinese \emph{'waen}) (KANJIDIC2; Middle Chinese in Baxter's transcription).
\end{cousinweb}

\begin{grammarlens}[title={What you've built}]
One more word for getting around.
\end{grammarlens}

\subsection*{Your turn}
\begin{itemize}
\raggedright
  \item \emph{Say it:} \emph{wan}
  \item \emph{Say it:} \emph{wan}, once more
  \item \emph{Say it:} \emph{wan no naka wa nami ga shizuka desu} — Inside the bay the waves are calm.
  \item \emph{Recall:} say the Japanese for to sparkle, then the Japanese for an island country, then say \emph{wan} again
\end{itemize}

\begin{tcolorbox}[breakable,colback=teal!4,colframe=teal!35!black,title={Before you move on}]
What does \ja{わん} mean? (\textquotedblleft{}A bay\textquotedblright{}.) Say it once more.
\end{tcolorbox}
\section[\texorpdfstring{\ja{かいきょう} (kaikyō)}{かいきょう (kaikyō)}]{\ja{かいきょう} (kaikyō) — a strait}

\label{lesson:JA-C383-kaikyou}



\begin{quote}
Before the new one: say the Japanese for an island country, then the Japanese for a bay.
\end{quote}

\subsection*{You'll want to know: \ja{かいきょう}}
\textbf{\ja{かいきょう}} — \emph{kaikyō} — \textquotedblleft{}a strait\textquotedblright{}.

\begin{quote}
\textbf{\ja{ふねで} \ja{かいきょうを} \ja{わたりました。}} — \emph{fune de kaikyō o watarimashita} — \textquotedblleft{}We crossed the strait by boat.\textquotedblright{}
\end{quote}

\begin{cousinweb}
\textbf{\ja{かいきょう}} (\emph{kaikyō}) is Sino-Japanese (kango): two Chinese characters read with their on'yomi, readings borrowed from Chinese. They are \emph{kai}, \textquotedblleft{}sea\textquotedblright{} (Middle Chinese \emph{xojX}) and \emph{kyō}, \textquotedblleft{}gorge\textquotedblright{} (\emph{heap}) (KANJIDIC2; Middle Chinese in Baxter's transcription).
\end{cousinweb}

\begin{grammarlens}[title={What you've built}]
One more word for getting around.
\end{grammarlens}

\subsection*{Your turn}
\begin{itemize}
\raggedright
  \item \emph{Say it:} \emph{kaikyō}
  \item \emph{Say it:} \emph{kaikyō}, once more
  \item \emph{Say it:} \emph{fune de kaikyō o watarimashita} — We crossed the strait by boat.
  \item \emph{Recall:} say the Japanese for an island country, then the Japanese for a bay, then say \emph{kaikyō} again
\end{itemize}

\begin{tcolorbox}[breakable,colback=teal!4,colframe=teal!35!black,title={Before you move on}]
What does \ja{かいきょう} mean? (\textquotedblleft{}A strait\textquotedblright{}.) Say it once more.
\end{tcolorbox}
\section[\texorpdfstring{\ja{かいすいよく} (kaisuiyoku)}{かいすいよく (kaisuiyoku)}]{\ja{かいすいよく} (kaisuiyoku) — swimming in the sea}

\label{lesson:JA-C383-kaisuiyoku}



\begin{quote}
Before the new one: say the Japanese for a bay, then the Japanese for a strait.
\end{quote}

\subsection*{You'll want to know: \ja{かいすいよく}}
\textbf{\ja{かいすいよく}} — \emph{kaisuiyoku} — \textquotedblleft{}swimming in the sea\textquotedblright{}.

\begin{quote}
\textbf{\ja{なつやすみに} \ja{かぞくで} \ja{かいすいよくに} \ja{いきました。}} — \emph{natsuyasumi ni kazoku de kaisuiyoku ni ikimashita} — \textquotedblleft{}In the summer holidays my family went swimming in the sea.\textquotedblright{}
\end{quote}

\begin{cousinweb}
\textbf{\ja{かいすいよく}} (\emph{kaisuiyoku}) is Sino-Japanese (kango): three Chinese characters read with their on'yomi, readings borrowed from Chinese. They are \emph{kai}, \textquotedblleft{}sea\textquotedblright{} (Middle Chinese \emph{xojX}), \emph{sui}, \textquotedblleft{}water\textquotedblright{} (\emph{sywijX}) and \emph{yoku}, \textquotedblleft{}bathe\textquotedblright{} (\emph{yowk}) (KANJIDIC2; Middle Chinese in Baxter's transcription).
\end{cousinweb}

\begin{grammarlens}[title={What you've built}]
One more word for getting around.
\end{grammarlens}

\subsection*{Your turn}
\begin{itemize}
\raggedright
  \item \emph{Say it:} \emph{kaisuiyoku}
  \item \emph{Say it:} \emph{kaisuiyoku}, once more
  \item \emph{Say it:} \emph{natsuyasumi ni kazoku de kaisuiyoku ni ikimashita} — In the summer holidays my family went swimming in the sea.
  \item \emph{Recall:} say the Japanese for a bay, then the Japanese for a strait, then say \emph{kaisuiyoku} again
\end{itemize}

\begin{tcolorbox}[breakable,colback=teal!4,colframe=teal!35!black,title={Before you move on}]
What does \ja{かいすいよく} mean? (\textquotedblleft{}Swimming in the sea\textquotedblright{}.) Say it once more.
\end{tcolorbox}
\section[\texorpdfstring{\ja{すきとおる} (sukitōru)}{すきとおる (sukitōru)}]{\ja{すきとおる} (sukitōru) — to be clear; to be see-through}

\label{lesson:JA-C383-sukitouru}



\begin{quote}
Before the new one: say the Japanese for a strait, then the Japanese for swimming in the sea.
\end{quote}

\subsection*{You'll want to know: \ja{すきとおる}}
\textbf{\ja{すきとおる}} — \emph{sukitōru} — \textquotedblleft{}to be clear; to be see-through\textquotedblright{}.

\begin{quote}
\textbf{\ja{この} \ja{うみは} \ja{みずが} \ja{すきとおっています。}} — \emph{kono umi wa mizu ga sukitōtte imasu} — \textquotedblleft{}The water of this sea is crystal clear.\textquotedblright{}
\end{quote}

Polite \textbf{\ja{すきとおります}}, \emph{sukitōrimasu}; a godan verb. Past (polite): \textbf{\ja{いけの} \ja{みずが} \ja{すきとおりました。}}, \emph{ike no mizu ga sukitōrimashita}.

\begin{cousinweb}
\textbf{\ja{すきとおる}} (\emph{sukitōru}) is a native Japanese word (wago), inherited rather than borrowed from Chinese; its kanji was attached later as a kun'yomi (JMdict).
\end{cousinweb}

\begin{grammarlens}[title={What you've built}]
One more word for getting around.
\end{grammarlens}

\subsection*{Your turn}
\begin{itemize}
\raggedright
  \item \emph{Say it:} \emph{sukitōru}
  \item \emph{Say it:} \emph{sukitōru}, once more
  \item \emph{Say it:} \emph{kono umi wa mizu ga sukitōtte imasu} — The water of this sea is crystal clear.
  \item \emph{Recall:} say the Japanese for a strait, then the Japanese for swimming in the sea, then say \emph{sukitōru} again
\end{itemize}

\begin{tcolorbox}[breakable,colback=teal!4,colframe=teal!35!black,title={Before you move on}]
What does \ja{すきとおる} mean? (\textquotedblleft{}To be clear; to be see-through\textquotedblright{}.) Say it once more.
\end{tcolorbox}
\section[(dialogue)]{Looking back: from A Crowded Carriage to An Island Country}

\label{lesson:JA-R383-review-378-383}



\begin{quote}
Say the Japanese for full (of people) and congestion.
\end{quote}

\subsection*{Your turn}
Say these aloud:

\begin{itemize}
\raggedright
  \item full (of people), congestion, packed tight, to cram in, to stand up
  \item to sit down, a railway carriage, an aisle, inside the train, a fare
  \item a station lunch box, a train conductor, to oversleep, dozing off, a mouthful
  \item a journey by boat, a sea route, a shore, to row, to drift
  \item the coast, a sandy beach, a cape, the horizon (over the sea), to sparkle
  \item an island country, a bay, a strait, swimming in the sea, to be see-through
\end{itemize}

\begin{tcolorbox}[breakable,colback=teal!4,colframe=teal!35!black,title={Before you move on}]
Full (of people)? (\textbf{\ja{まんいん}}, \emph{man'in}.) Congestion? (\textbf{\ja{こんざつ}}, \emph{konzatsu}.) Packed tight? (\textbf{\ja{ぎゅうぎゅう}}, \emph{gyūgyū}.) To cram in? (\textbf{\ja{つめこむ}}, \emph{tsumekomu}.) To stand up? (\textbf{\ja{たちあがる}}, \emph{tachiagaru}.) To sit down? (\textbf{\ja{すわる}}, \emph{suwaru}.) A railway carriage? (\textbf{\ja{しゃりょう}}, \emph{sharyō}.) An aisle? (\textbf{\ja{つうろ}}, \emph{tsūro}.) Inside the train? (\textbf{\ja{しゃない}}, \emph{shanai}.) A fare? (\textbf{\ja{うんちん}}, \emph{unchin}.) A station lunch box? (\textbf{\ja{えきべん}}, \emph{ekiben}.) A train conductor? (\textbf{\ja{しゃしょう}}, \emph{shashō}.) To oversleep? (\textbf{\ja{ねすごす}}, \emph{nesugosu}.) Dozing off? (\textbf{\ja{うとうと}}, \emph{utōto}.) A mouthful? (\textbf{\ja{ひとくち}}, \emph{hitokuchi}.) A journey by boat? (\textbf{\ja{ふなたび}}, \emph{funatabi}.) A sea route? (\textbf{\ja{こうろ}}, \emph{kōro}.) A shore? (\textbf{\ja{きし}}, \emph{kishi}.) To row? (\textbf{\ja{こぐ}}, \emph{kogu}.) To drift? (\textbf{\ja{ただよう}}, \emph{tadayou}.) The coast? (\textbf{\ja{かいがん}}, \emph{kaigan}.) A sandy beach? (\textbf{\ja{すなはま}}, \emph{sunahama}.) A cape? (\textbf{\ja{みさき}}, \emph{misaki}.) The horizon (over the sea)? (\textbf{\ja{すいへいせん}}, \emph{suiheisen}.) To sparkle? (\textbf{\ja{かがやく}}, \emph{kagayaku}.) An island country? (\textbf{\ja{しまぐに}}, \emph{shimaguni}.) A bay? (\textbf{\ja{わん}}, \emph{wan}.) A strait? (\textbf{\ja{かいきょう}}, \emph{kaikyō}.) Swimming in the sea? (\textbf{\ja{かいすいよく}}, \emph{kaisuiyoku}.) To be see-through? (\textbf{\ja{すきとおる}}, \emph{sukitōru}.)
\end{tcolorbox}
